\section{Defect Classes and Reactions}\label{sec:classes}

The primary irradiation defects in hcp zirconium are vacancies, self-interstitial
atoms and their clusters. Because the lattice is elastically and
crystallographically anisotropic, a cluster is not specified by its size alone:
its Burgers-vector character and habit plane determine its contribution to
irradiation growth, to hardening, and to the sink field seen by every other
species. The observed loop population in irradiated Zr and its alloys is
conventionally grouped into \aloop{} loops and \cloop{}-component
loops~\citep{kelly1973characterization,ChoiKim2013,tournadre2012experimental},
with
\begin{equation}
  \bvec_a=\tfrac13\langle11\bar20\rangle,
  \qquad
  \bvec_c \parallel [0001],
  \label{eq:burgers-general}
\end{equation}
the latter carrying a Burgers-vector component along the hexagonal axis.

Four principal cluster families span that population, together with the two free
monomers~\citep{Carpenter1988,Northwood1977,OnimusBechade2012,Yi2013,Doriot2015}:
\begin{equation}
  I_n^{a},\qquad V_n^{a},\qquad S_n,\qquad V_n^{c},
  \label{eq:zr-loop-classes}
\end{equation}
where $I_n^{a}$ is an interstitial prismatic \aloop{} loop of $n$ excess
self-interstitials on a $\{10\bar10\}$ habit, $V_n^{a}$ a vacancy prismatic loop
of the same habit and Burgers vector but opposite character, $S_n$ a basal
stacking-fault pyramid of $n$ aggregated vacancies bounded by basal
$\tfrac16\langle20\bar23\rangle$ and prismatic facets and formed directly from
collapsed cascade vacancy debris, and $V_n^{c}$ a basal vacancy
\cloop{}-component loop, which nucleates above a temperature- and
dose-rate-dependent threshold, in part by unfaulting of supercritical pyramids,
and which drives breakaway irradiation growth at high
fluence~\citep{holt1986c,HoltCausey1993,Yi2013}. The free monomers are the
mono-interstitial $C_i$ and the monovacancy $C_v$. Neither helium nor a
three-dimensional void population is carried, because Zr alloys generate a
negligible $(n,\alpha)$ helium yield and multi-vacancy cascade debris in
$\alpha$-Zr collapses onto basal or prismatic habits rather than forming compact
cavities~\citep{Northwood1977,OnimusBechade2012}.

This classification separates character, interstitial or vacancy, from
crystallography, prismatic or basal, and both distinctions carry weight. Two
families sharing a habit plane and a Burgers vector but differing in character,
$I_n^a$ and $V_n^a$, respond to the same mobile fluxes with opposite sign, while
two families sharing a character but differing in habit, $V_n^a$ and $V_n^c$,
present differently oriented capture cross-sections to an anisotropic diffusion
field and therefore compete for the same vacancies on unequal terms. It is that
second competition, taken up in \cref{sec:coexistence}, which the self-consistent
formulation exists to describe without a free parameter.

\Cref{eq:zr-loop-classes} is size-resolved, each family being an infinite ladder
in $n$. The model solved here replaces the ladder by a small number of moments of
the size distribution per family (\cref{sec:moments}) and resolves instead the
crystallographic variants that a size-resolved but variant-lumped model cannot
distinguish. Writing the character as $s\in\{i,v\}$ and the habit as
$k\in\{a_1,a_2,a_3,c_0,c_f,c_p\}$, the families carried are
\begin{equation}
  \mathcal F=
  \underbrace{\{(i,a_1),(i,a_2),(i,a_3)\}}_{\text{prismatic interstitial}}
  \cup
  \underbrace{\{(v,a_1),(v,a_2),(v,a_3)\}}_{\text{prismatic vacancy}}
  \cup
  \underbrace{\{(v,c_0),(v,c_f),(v,c_p)\}}_{\text{basal vacancy chain}},
  \label{eq:families}
\end{equation}
nine in all, of which eight are dislocation loops and one is not: $c_0$ is the
stacking-fault pyramid, the precursor from which the basal loops form. We write
$\mathcal F_L=\mathcal F\setminus\{(v,c_0)\}$ for the eight loop families, and
every statement below invoking a perimeter, a line density or a coalescence
channel is a statement about $\mathcal F_L$.

Three features distinguish \cref{eq:families} from the lumped, stress-aligned
pair of the earlier mean-field model, and each is a statement about
crystallography rather than about the solver. The prismatic loops are resolved
into their three variants of both characters, so that each carries its own habit
normal $\hat{\bm a}_k$, the resolved stress on each is computed rather than
parameterized, and the response to an arbitrary stress state is a redistribution
among the three rather than a shift of one scalar; at zero deviatoric stress the
three are equivalent by symmetry, receive equal shares, and sum to the lumped
population exactly, which is the regression test on the change. There is no basal
interstitial family, because basal interstitial \cloop{} loops are not a reported
population in neutron-irradiated Zr; a formal $c_i$ population is retained in the
growth inequalities of \cref{sec:coexistence} for symmetry but is not a state
variable. And the basal vacancy population is a three-state chain
$c_0\to c_f\to c_p$, in which the pyramid is not planar and stores its vacancies
with a nearly isotropic relaxation volume while the faulted and perfect states
are planar and store different vacancy content per unit area; both transitions
are thermally activated and both are one-way.

% The hcp cell and every Burgers vector the model carries. Drawn rather than
% included as a raster, for three reasons: the raster was drawn at an arbitrary
% axial ratio where alpha-Zr's is 1.5944 and the distinction matters here (it is
% what separates the two basal states); it carried no Burgers vectors, which are
% the thing a reader needs from this figure; and its colors were unrelated to
% the family colors every other figure in the paper uses.
%
% PROJECTION: elevation 20 deg, azimuth -45 deg, applied through TikZ's own
% x/y/z keys, so the coordinates below are the crystal's and the axial ratio is
% exact by construction rather than by drafting. The azimuth is -45 and not -60
% for a reason: at -60 the hexagon's vertices project in coincident pairs, the
% left and right edges come out vertical, and the prism reads as a rectangular
% block rather than a hexagonal one.
%
% The three panels sit in bottom-aligned minipages so that the (a)/(b)/(c)
% labels land on one line; aligning the tikzpictures themselves by `baseline`
% cannot do that, because each panel's origin sits at a different height within
% its own drawing.
\definecolor{famA1}{HTML}{C62828}
\definecolor{famA2}{HTML}{2E7D32}
\definecolor{famA3}{HTML}{B38600}
\definecolor{famA1v}{HTML}{00838F}
\definecolor{famA2v}{HTML}{00695C}
\definecolor{famA3v}{HTML}{6A1B9A}
\definecolor{famCf}{HTML}{1F4FBF}
\definecolor{famCp}{HTML}{4527A0}

\begin{figure}[htbp]
\centering
%---------------------------------------------------------------- (a) the cell
\begin{minipage}[b]{0.30\textwidth}\centering
\begin{tikzpicture}[
    x={(1.20cm,-0.411cm)}, y={(1.20cm,0.411cm)}, z={(0cm,1.60cm)},
    >=Latex, line join=round]
  \def\h{1.5944}                      % c/a for alpha-Zr, not sqrt(8/3)
  \foreach \i/\X/\Y in {0/1.000/0.000, 1/0.500/0.866, 2/-0.500/0.866,
                        3/-1.000/0.000, 4/-0.500/-0.866, 5/0.500/-0.866} {
    \coordinate (B\i) at (\X,\Y,0);
    \coordinate (T\i) at (\X,\Y,\h);
  }

  % --- what the camera cannot see, first and dashed
  \draw[gray!45, dashed] (B1) -- (B2) -- (B3) -- (B4);
  \draw[gray!45, dashed] (B2) -- (T2);

  % --- the second-order {11-20} plane: a diagonal THROUGH the cell, normal to
  %     an <a> direction. Lightly filled: solid, it hides the cell it cuts.
  \draw[famA3, densely dashed, fill=famA3!14]
    (0,0.866,0) -- (0,-0.866,0) -- (0,-0.866,\h) -- (0,0.866,\h) -- cycle;

  % --- a first-order {10-10} face, front-facing at this azimuth
  \fill[famA1!28, draw=famA1!80]
    (B5) -- (B0) -- (T0) -- (T5) -- cycle;

  % --- the basal habit (0001). The TOP face is the visible one from here.
  \fill[famCf!25, draw=famCf!80]
    (T0) -- (T1) -- (T2) -- (T3) -- (T4) -- (T5) -- cycle;

  % --- the remaining visible edges
  \draw[gray!70] (B4) -- (B5) -- (B0) -- (B1);
  \draw[gray!70] (B0) -- (T0);  \draw[gray!70] (B1) -- (T1);
  \draw[gray!70] (B3) -- (T3);  \draw[gray!70] (B4) -- (T4);
  \draw[gray!70] (B5) -- (T5);

  % --- the two shortest lattice translations
  \draw[->, very thick, famA1] (0,0,\h) -- (1.0,0,\h);
  \node[famA1, font=\footnotesize] at (1.36,0.12,\h) {$a$};
  \draw[->, very thick, famCp] (0,0,0) -- (0,0,\h);
  \node[famCp, font=\footnotesize] at (0.20,0.34,{\h*0.50}) {$c$};
\end{tikzpicture}
\\[2pt]
{\scriptsize\color{gray}$c/a=1.5944$}\\[1pt]
{\small (a)}
\end{minipage}
\hfill
%------------------------------------------------- (b) plan view, prismatic set
\begin{minipage}[b]{0.32\textwidth}\centering
\begin{tikzpicture}[scale=0.98, >=Latex, line join=round]
  % the {10-10} first-order traces bisect the <a> directions
  \foreach \A in {30,90,150}
    \draw[famA3!55, densely dashed] (\A-180:1.28) -- (\A:1.28);
  \node[font=\scriptsize, famA3, anchor=south west] at (30:1.22)
       {$\{10\bar10\}$};

  % basal section, viewed down [0001]
  \draw[gray!60, fill=gray!6]
    (0:1) -- (60:1) -- (120:1) -- (180:1) -- (240:1) -- (300:1) -- cycle;

  % the three prismatic Burgers vectors, interstitial (long, solid) and vacancy
  % (short, dashed): the two characters share habit and b and differ only in
  % what they store, which no arrow can show.
  \foreach \A/\ci/\cv/\lbl/\pos in {%
      0/famA1/famA1v/{$a_1,a_1^{v}$}/{right},
      60/famA2/famA2v/{$a_2,a_2^{v}$}/{above right},
      120/famA3/famA3v/{$a_3,a_3^{v}$}/{above left}} {
    \draw[->, very thick, \ci] (0,0) -- (\A:1.00);
    \draw[->, thick, \cv, densely dashed] (0,0) -- (\A:0.58);
    \node[font=\scriptsize, \ci, \pos] at (\A:1.04) {\lbl};
  }
  \node[font=\scriptsize, anchor=north] at (0,-1.32) {$[0001]$ into the page};
\end{tikzpicture}
\\[2pt]
{\scriptsize $\bm b=\tfrac13\langle11\bar20\rangle$ for all six\\[1pt]
 solid: interstitial\quad dashed: vacancy}\\[1pt]
{\small (b)}
\end{minipage}
\hfill
%----------------------------------------------------- (c) elevation, basal set
\begin{minipage}[b]{0.32\textwidth}\centering
% ONLY the family names sit beside the arrows. The Burgers vectors go in the
% legend below: at this scale a symbol long enough to name b crowds either the
% c-axis scale bar or the other family's arrow, whichever side it is put on.
\begin{tikzpicture}[scale=0.90, >=Latex, line join=round]
  \def\h{1.5944}
  % the c axis, for scale, clear of everything
  \draw[gray!55, |<->|] (1.30,0) -- (1.30,\h);
  \node[gray!70, font=\scriptsize, anchor=west] at (1.35,{\h/2}) {$c$};

  % the (0001) habit, seen edge-on
  \draw[famCf!70, fill=famCf!16] (-1.25,-0.05) rectangle (1.10,0.05);
  \node[font=\scriptsize, famCf, anchor=north] at (0.40,-0.09) {$(0001)$};

  % c_p: perfect, b = [0001], edge component the full c
  \draw[->, very thick, famCp] (0.62,0.05) -- (0.62,\h);
  \node[font=\scriptsize, famCp, anchor=south] at (0.62,{\h+0.05}) {$c_p$};

  % c_f: faulted, b = 1/6<20-23>: inclined, edge component c/2
  \draw[->, very thick, famCf] (-0.15,0.05) -- (-1.00,{\h/2});
  \draw[famCf!50, densely dotted] (-1.00,{\h/2}) -- (-0.15,{\h/2});
  \draw[famCf!50, densely dotted] (-0.15,0.05) -- (-0.15,{\h/2});
  \node[font=\scriptsize, famCf, anchor=south] at (-0.60,{\h/2+0.12})
       {$c_f$};
  \node[font=\scriptsize, famCf, anchor=west] at (-0.10,{\h/4}) {$c/2$};
\end{tikzpicture}
\\[2pt]
{\scriptsize
 \textcolor{famCf}{$c_f$}: $\bm b=\tfrac16\langle20\bar23\rangle$,
 edge component $c/2$\\[1pt]
 \textcolor{famCp}{$c_p$}: $\bm b=[0001]$, edge component $c$}\\[1pt]
{\small (c)}
\end{minipage}

\caption{The hcp cell of $\alpha$-Zr and every Burgers vector the model carries,
drawn at the physical axial ratio $c/a=1.5944$ rather than the ideal
$\sqrt{8/3}=1.6330$. \textbf{(a)}~The cell, with the three planes that matter:
the basal $(0001)$ plane (blue), the habit of $c_0$, $c_f$ and $c_p$; a
first-order prismatic $\{10\bar10\}$ face (red), a side face of the prism; and a
second-order $\{11\bar20\}$ plane (yellow, dashed), a diagonal \emph{through}
the cell rather than a face of it. The arrows are the two shortest lattice
translations, $a=\tfrac13\langle11\bar20\rangle$ in the basal plane and
$c=[0001]$ along the cell axis. \textbf{(b)}~Plan view down $[0001]$: the three
prismatic variants at $60^\circ$ to one another, each drawn twice --- long and
solid for the interstitial family, short and dashed for its vacancy partner,
which shares the habit plane and the Burgers vector and differs only in what it
stores. The dashed yellow lines are the $\{10\bar10\}$ traces, which bisect the
Burgers directions. \textbf{(c)}~The basal pair, seen edge-on. The perfect state
$c_p$ has $\bm b=[0001]$ normal to its habit, so its edge component is the full
$c$; the faulted state $c_f$ has the inclined
$\bm b=\tfrac16\langle20\bar23\rangle$, whose edge component is $c/2$ --- one
removed basal layer against two. That factor of two is why
$\lambda_{c_p}=\lambda_{c_f}/\sqrt2$ in \cref{tab:cd-families}, and it is why
the axial ratio has to be drawn correctly. The pyramid $c_0$ has no Burgers
vector and so appears in neither panel. The inclined pyramidal planes containing
$\langle c+a\rangle=\tfrac13\langle11\bar23\rangle$ carry \emph{no} family: that
vector is the glide vector accommodating $c$-axis strain in deformation, and its
planes are not observed habit planes for irradiation-induced loops, so the
absence is an experimental statement and not a truncation.}
\label{fig:hcp-cell}
\end{figure}


The two habits are mutually orthogonal: the normals $\hat{\bm a}_k$ of the three
prismatic variants lie in $(0001)$ while the basal normal is $[0001]$ itself.
That orthogonality is why no single electron-beam direction images both
populations at true size, a zone axis showing one family face-on viewing the
other edge-on, so that any comparison of \cref{eq:families} against
transmission-electron-microscope densities must combine at least two zone axes
and a single-axis micrograph constrains at most one habit. It is also why the
calibration of \cref{sec:calibration} treats the \aloop{} and \cloop{} data as
separate target sets rather than as one.

Each family carries a character sign $\eta_s$, $+1$ for interstitial and $-1$ for
vacancy, and three spatially resolved fields: the cluster density $n^k_{sL}(\xv)$
in clusters per host atom, the defect content $c^k_{sL}(\xv)$ in point defects
stored per host atom, and the second content $q^k_{sL}(\xv)$ of
\cref{sec:moments}, which carries the width of the size distribution. All three
are dimensionless, so a number density per unit volume is $n/\Omega$ and a
stored-defect density is $c/\Omega$, with $\Omega$ the atomic volume. The
quantity converting loop area into stored defects is not $|\bvec|$ but the edge
component $\bvec\!\cdot\!\nvec$, since a loop of radius $R$ holds
$\pi R^{2}(\bvec\!\cdot\!\nvec)/\Omega$ defects, so that the mean radius, the
mean size and the line density of family $(s,k)$ follow from its content and
density as
\begin{equation}
  R^k_{sL}=\lambda_k\sqrt{\frac{c^k_{sL}}{n^k_{sL}}},
  \qquad
  \lambda_k=\sqrt{\frac{\Omega}{\pi\,(\bvec\!\cdot\!\nvec)_k}},
  \qquad
  m^k_{sL}=\frac{c^k_{sL}}{n^k_{sL}},
  \qquad
  \rho^k_{sL}=2\pi R^k_{sL}\,n^k_{sL},
  \label{eq:radius}
\end{equation}
with one geometric length per habit: $\lambda_a$ shared by all six prismatic
families, and $\lambda_{c_f}$, $\lambda_{c_p}$ for the two basal states.
\Cref{eq:radius} describes a planar object and does not apply to the pyramid,
whose stored vacancies fill a volume rather than bound an area; its size follows
instead from the volumetric relation
\begin{equation}
  R^{c_0}_{vL}=\left(\frac{m^{c_0}_{vL}\,\Omega}{\sqrt8}\right)^{1/3},
  \label{eq:Rsfp}
\end{equation}
so that the pyramid carries no $\lambda_k$, no Burgers vector and no line
density. What it does carry is a capture cross-section, and \cref{sec:formulation}
shows that a cross-section is the only thing either the growth law or the sink
strengths ever ask of a family. \Cref{tab:cd-families} lists the nine families
with the crystallography that fixes their geometry.

\begin{table}[htbp]
\caption{The nine families of \cref{eq:families}, with the crystallography that
fixes each one's geometry. $\lambda_k=\sqrt{\Omega/\pi(\bvec\!\cdot\!\nvec)_k}$
is the length converting $\sqrt{m}$ into a radius, \cref{eq:radius}.
The lattice constants are $a=\SI{0.323}{\nano\meter}$,
$c=\SI{0.515}{\nano\meter}$, hence $c/a=1.5944$ --- the physical ratio and not
the ideal $\sqrt{8/3}=1.6330$ --- and
$\Omega=(\sqrt3/4)a^2c=\SI{2.3266e-29}{\cubic\meter}$.}
\label{tab:cd-families}
\centering\small
\begin{tabularx}{\textwidth}{@{}llccrX@{}}
\toprule
Family & Habit & $\bvec$ & Character & $\lambda_k$ (nm) & Comment \\
\midrule
$a_1,a_2,a_3$ & $\{10\bar10\}$ & $\tfrac13\langle11\bar20\rangle$ & $i$ &
  0.151 & The dominant early population; $\bvec\!\cdot\!\nvec=a$. \\
$a_1,a_2,a_3$ & $\{10\bar10\}$ & $\tfrac13\langle11\bar20\rangle$ & $v$ &
  0.151 & Same habit and $|\bvec|$ as the interstitial variants; separated from
  them only by capture character. \\
$c_f$ & $(0001)$ & $\tfrac16\langle20\bar23\rangle$ & $v$ &
  0.170 & Faulted, $\bvec\!\cdot\!\nvec=c/2$, one removed basal layer. The
  dominant observed \cloop{} loop. \\
$c_p$ & $(0001)$ & $[0001]$ & $v$ &
  0.120 & Perfect, $\bvec\!\cdot\!\nvec=c$, two removed layers, hence
  $\lambda_{c_p}=\lambda_{c_f}/\sqrt2$. \\
$c_0$ & $(0001)$ base & partials & $v$ &
  --- & The stacking-fault pyramid: not planar, so \cref{eq:radius} is replaced
  by \cref{eq:Rsfp} and it carries no $\lambda_k$. Precursor to $c_f$. \\
\bottomrule
\end{tabularx}
\end{table}

The coupled kinetics are generated from a directed multigraph, the reaction
admissibility graph, whose vertices are the admissible cluster states and whose
edges are the elementary reactions permitted by the polarity, habit-plane and
mobility constraints above. The master equations follow by a deterministic graph
walk, the time derivative at each vertex being the signed sum of the rate kernels
of all incoming and outgoing edges incident on it. Vertices partition by
polarity,
\begin{align}
  \mathcal V^{+} &= \{C_i,C_{2i},C_{3i}\}\cup\{(i,a_k)\}, \\
  \mathcal V^{-} &= \{C_v\}\cup\{(v,a_k)\}\cup\{(v,c_0),(v,c_f),(v,c_p)\},
\end{align}
so that the interstitial prismatic families lie on the positive branch and the
six vacancy families on the negative one. \Cref{tab:rag-processes} catalogs the
fourteen edge classes and \cref{fig:rag} draws the graph, and both are drawn
against the equations summarized in \cref{sec:summary-equations} rather than
against a description of them.

Three of those classes deserve comment, because they are absent from the
mean-field treatment this model replaces and are easily overlooked. Cascade
production is not one edge but two partitions of the same displacement rate: a
part $G_m$ deposited into the mobile species, and a part $J_k$ that nucleates
loops directly, with its own share for the prismatic families of \emph{both}
characters as well as for the basal family. Alongside that cascade route the
interstitial prismatic family has a second, homogeneous nucleation channel fed by
$2i+2i$ and $i+3i$ reactions among the mobile clusters, so that the mobile ladder
and the loop populations are not separable even in principle. And the grain
boundary appears twice, as a Dirichlet condition on the mobile species and as a
size-gated removal acting on the loops themselves; the second has no counterpart
in a mean-field treatment, where a loop has no size relative to a boundary that
is itself only a sink strength.

\begin{table}[htbp]
\caption{Admissible process classes of the irradiated-Zr reaction admissibility
graph. Each elementary reaction maps to one abstract edge class and one rate
kernel; the explicit kernels appear in \cref{sec:formulation} at the equation
indicated. P14 has no counterpart in the mean-field treatment and is introduced
in \cref{sec:gb-loops}.}
\label{tab:rag-processes}
\centering\small
\begin{tabularx}{\textwidth}{@{}l>{\raggedright\arraybackslash}p{0.32\textwidth}Xl@{}}
\toprule
& Process & Kernel & Given in \\
\midrule
P1  & vacancy--interstitial recombination, and vacancy capture by a mobile
      cluster & $R_{i,v}$, $R_{v,2i}$, $R_{v,3i}$ & \cref{eq:rates} \\
P2  & clustering among mobile species & $\beta^{m,n}$ & \cref{eq:rates} \\
P3  & thermal dissociation of a mobile cluster & $\alpha^{m,n}$ &
      \cref{eq:rates} \\
P4  & cascade production into the mobile species & $G_m$ & \cref{eq:master-c} \\
P5  & cascade nucleation of loops, both characters & $J_k$ & \cref{eq:block-n} \\
P6  & homogeneous loop nucleation from mobile clusters &
      $R_{i,3i}+R_{2i,2i}$ & \cref{eq:block-n} \\
P7  & capture at a loop or the pyramid, of either polarity & $\phi_{km}$ &
      \cref{eq:growth} \\
P8  & peripheral emission from a vacancy loop & $S_{vL}$ & \cref{eq:Semit} \\
P9  & loop--loop coalescence & $\varphi_{LL}$ & \cref{eq:avrami} \\
P10 & loop--network coalescence & $\varphi_{LN}$ & \cref{eq:avrami} \\
P11 & absorption of a mobile species at the network & $R_{m,s}$ &
      \cref{eq:sinkconc} \\
P12 & basal chain transfer, $c_0\to c_f\to c_p$ &
      $\nu_{\mathrm{col}},\nu_{\mathrm{uf}}$ & \cref{eq:basal-chain} \\
P13 & absorption of a mobile species at a grain boundary & Dirichlet &
      \cref{eq:bc} \\
P14 & size-gated absorption of \emph{loops} at a grain boundary & $\nu_{gb}$ &
      \cref{eq:gbloop} \\
\bottomrule
\end{tabularx}
\end{table}

% The reaction admissibility graph for the nine families. Drawn here rather
% than lifted from the mean-field treatment: that graph is a size-resolved
% ladder over four families, and this model carries moments over nine, with a
% basal chain and a grain-boundary loop channel the ladder has no vertex or
% edge for.
%
% EVERY EDGE IS CHECKED AGAINST THE SOLVER, not against the prose. The earlier
% version of this figure omitted the cascade sources of C_2i, C_3i and of the
% prismatic vacancy loops; the mobile clustering ladder and its dissociation
% edges entirely; the homogeneous loop-nucleation channel that 2i+2i and i+3i
% feed; and the capture edges onto the basal family. Each of those appears in
% `rate_equations_core.h`, and each is now drawn.
%
% The four mobile species are boxed together so that capture onto a loop can be
% ONE arrow per family rather than four: the growth law sums over all four
% species with a sign set by whether the polarity matches the family's, so the
% bundle is faithful and the alternative is twelve crossing arrows.
\begin{figure}[htbp]
\centering
\begin{tikzpicture}[
  font=\small, >=Latex,
  mob/.style={draw, rounded corners=2pt, minimum height=7mm,
              minimum width=10mm, fill=blue!8, draw=blue!55},
  ipop/.style={draw, rounded corners=2pt, align=center, fill=red!8,
               draw=red!60, minimum height=8mm, inner sep=3pt},
  vpop/.style={draw, rounded corners=2pt, align=center, fill=cyan!10,
               draw=cyan!55!black, minimum height=8mm, inner sep=3pt},
  bpop/.style={draw, rounded corners=2pt, align=center, fill=violet!10,
               draw=violet!60, minimum height=8mm, inner sep=3pt},
  ext/.style={draw, align=center, fill=black!6, draw=black!45,
              minimum height=7mm, inner sep=3pt},
  clus/.style={->, thick, blue!70},
  capture/.style={->, very thick, blue!60},
  diss/.style={->, dashed, orange!85!black},
  rec/.style={<->, very thick, red!75},
  coal/.style={->, thick, orange!80!black},
  chain/.style={->, very thick, violet!70},
  snk/.style={->, thick, black!50},
  gbl/.style={->, very thick, green!45!black},
  src/.style={->, thick, green!55!black}]

  % ---- cascade source ---------------------------------------------------
  \node[ext] (G) at (-1.2,5.5) {cascade source $G$};

  % ---- the mobile block -------------------------------------------------
  \node[mob] (Cv)  at (-6.2,3.0) {$C_v$};
  \node[mob] (Ci)  at (-2.6,3.0) {$C_i$};
  \node[mob] (C2i) at (-0.2,3.0) {$C_{2i}$};
  \node[mob] (C3i) at ( 2.2,3.0) {$C_{3i}$};
  \node[draw, rounded corners=3pt, dashed, draw=blue!45, inner sep=6pt,
        fit=(Cv)(Ci)(C2i)(C3i), label={[font=\scriptsize, text=blue!60,
        anchor=south east]south east:{mobile species}}] (MOB) {};

  % ---- immobile families ------------------------------------------------
  \node[ipop] (ai) at (-3.4,0.3)
       {$(i,a_1),(i,a_2),(i,a_3)$\\[-1pt]
        \scriptsize prismatic $\langle a\rangle$, interstitial};
  \node[vpop] (av) at (-3.4,-2.2)
       {$(v,a_1),(v,a_2),(v,a_3)$\\[-1pt]
        \scriptsize prismatic $\langle a\rangle$, vacancy};
  \node[bpop] (c0) at (-5.6,-4.7) {$c_0$\\[-1pt]\scriptsize SFP};
  \node[bpop] (cf) at (-2.6,-4.7) {$c_f$\\[-1pt]\scriptsize faulted};
  \node[bpop] (cp) at ( 0.4,-4.7) {$c_p$\\[-1pt]\scriptsize perfect};

  % ---- external reservoirs ----------------------------------------------
  \node[ext] (NET) at (5.4,0.3)  {network\\[-1pt]\scriptsize $\rho_N$};
  \node[ext] (GB)  at (5.4,-3.0) {grain\\[-1pt]boundary};

  % ---- P4/P5 cascade sources --------------------------------------------
  \draw[src] (G.west) to[bend right=25]
       node[above left, pos=0.75, font=\scriptsize] {$G_m$ (P4)} (Cv.north);
  \draw[src] (G.south) -- node[right, pos=0.6, font=\scriptsize] {} (Ci.north);
  \draw[src] (G.south) -- (C2i.north);
  \draw[src] (G.south east) to[bend left=10] (C3i.north);
  \draw[src] (G.east) to[bend left=44]
       node[right, pos=0.35, font=\scriptsize] {$J_k$ (P5)} (ai.east);
  \draw[src] (G.east) to[bend left=52] (av.east);
  \draw[src, dashed] (G.east) to[bend left=58] (cf.east);

  % ---- P2 clustering, P3 dissociation ------------------------------------
  \draw[clus] (Ci.north east) to[bend left=30]
       node[above, font=\scriptsize] {$\beta^{i,i}$ (P2)} (C2i.north west);
  \draw[clus] (C2i.north east) to[bend left=30]
       node[above, font=\scriptsize] {$\beta^{i,2i}$} (C3i.north west);
  \draw[diss] (C2i.south west) to[bend left=30]
       node[below, font=\scriptsize] {$\alpha^{i,2i}$ (P3)} (Ci.south east);
  \draw[diss] (C3i.south west) to[bend left=30]
       node[below, font=\scriptsize] {$\alpha^{i,3i}$} (C2i.south east);

  % ---- P1 recombination and vacancy capture by clusters ------------------
  \draw[rec] (Cv) -- node[above, font=\scriptsize] {$R_{i,v}$ (P1)} (Ci);
  \draw[->, thick, red!60] (Cv.north east) to[bend left=22]
       node[above, pos=0.55, font=\scriptsize] {$R_{v,2i},R_{v,3i}$} (C3i.north west);

  % ---- P6 homogeneous loop nucleation -------------------------------------
  \draw[src, thick] (C3i.south) to[bend left=14]
       node[right, pos=0.55, font=\scriptsize]
       {$R_{i,3i}+R_{2i,2i}$ (P6)} (ai.north east);

  % ---- P7 capture at a loop or the SFP ------------------------------------
  \draw[capture] (MOB.south) to[bend right=6]
       node[left, pos=0.6, font=\scriptsize] {$\phi_{km}$ (P7)} (ai.north);
  \draw[capture] (MOB.south west) to[bend right=16] (av.north west);
  \draw[capture] (MOB.south west) to[bend right=34] (cf.north west);
  \draw[capture] (Cv.south) to[bend right=24] (c0.west);

  % ---- P8 peripheral emission ---------------------------------------------
  \draw[diss] (av.west) to[bend right=22]
       node[left, pos=0.5, font=\scriptsize] {$S_{vL}$ (P8)} (Cv.south west);
  \draw[diss] (cf.south) to[bend right=40] (Cv.south);

  % ---- P9 loop-loop coalescence -------------------------------------------
  \draw[coal] (ai.north west) to[out=150, in=210, looseness=5]
       node[left, font=\scriptsize] {$\varphi_{LL}$ (P9)} (ai.south west);
  \draw[coal] (cf.south west) to[out=-140, in=-40, looseness=5]
       node[below, font=\scriptsize] {$\varphi_{LL}$} (cf.south east);

  % ---- P10 loop-network, P11 mobile-network -------------------------------
  \draw[chain] (ai.east) -- node[above, font=\scriptsize]
       {$\varphi_{LN}$ (P10)} (NET.west);
  \draw[snk] (C3i.east) to[bend left=16]
       node[above, pos=0.7, font=\scriptsize] {$R_{m,s}$ (P11)} (NET.north);

  % ---- P12 basal chain ----------------------------------------------------
  \draw[chain] (c0) -- node[above, font=\scriptsize]
       {$\nu_{\mathrm{col}}$ (P12)} (cf);
  \draw[chain] (cf) -- node[above, font=\scriptsize]
       {$\nu_{\mathrm{uf}}$} (cp);

  % ---- P13 Dirichlet, P14 size-gated loop absorption ----------------------
  \draw[snk] (MOB.east) to[bend left=8]
       node[above, pos=0.6, font=\scriptsize] {Dirichlet (P13)} (GB.north);
  \draw[gbl] (av.east) to[bend right=6]
       node[above, pos=0.6, font=\scriptsize] {$\nu_{gb}$ (P14)} (GB.west);
  \draw[gbl] (cp.east) to[bend right=12] (GB.south);
\end{tikzpicture}
\caption{Reaction admissibility graph for the nine families of
\cref{eq:families}, with every edge class of \cref{tab:rag-processes}. The four
mobile species are boxed together because capture onto a loop sums over all of
them (\cref{eq:growth}), with a sign set by whether the species' polarity
matches the family's, so one bundled arrow per family is faithful where four
would be unreadable. Cascade production (green) feeds all four mobile species
and, as a separate partition of the same displacement rate, nucleates the
prismatic loops of both characters and the basal family directly; the dashed
source onto $c_f$ becomes a source onto $c_0$ instead when the basal chain is
active, which it is not in \cref{sec:results}. The interstitial prismatic
family has a second, homogeneous nucleation channel (P6) fed by
$2i+2i$ and $i+3i$ reactions among the mobile clusters, which is why the
mobile ladder and the loop populations are not separable. Dashed orange edges
run backwards from a population to the mobile pool: dissociation of a small
cluster (P3) and peripheral emission from a vacancy loop (P8). The three
prismatic variants of each character are drawn as one vertex for legibility;
they are separate populations that exchange nothing with one another and differ
only in habit normal, and the pyramid $c_0$ is not a loop --- it has no Burgers
vector, no perimeter and no coalescence channel, and enters only through a
capture cross-section.}
\label{fig:rag}
\end{figure}

